% SCA: exact original relation defect and the complete same-class allowance.
% Additive source. No preceding mathematical provider is changed.
\section{The literal relation defect and the original same-class action allowance}
\label{sec:SCA}

This calculation retains the original source, quotient action and observation.
It first evaluates the relation-valued term in AMC25, including its pairing
with the minimum section. It then evaluates the full original OMR19/OER13
action allowance to its already determined quadratic order. The observation
and residual terms are transported before their matching determinant terms
are cancelled. In particular a reference boundary in that allowance is not
held fixed while changing the identical boundary in its other summand.

\subsection{Original sources, primary coordinates and minimum lifts}

Fix the complete hypothetical packet of SXR1 and GQB1,
\[
 \begin{gathered}
 \rho=\tfrac12+\delta+i\gamma,
 \quad0<\delta<\tfrac12,\quad\gamma>2,\quad m\ge1,\\
 g=2\xi,\quad
 h(s)=\prod_{z\in\{\rho,\bar\rho,1-\rho,1-\bar\rho\}}(s-z)^m,
 \quad v_h=g/h,\quad
 w_h(y)=|v_h(1/2+iy)|^2/(2\pi),\\
 k=4l+1\ge5,\quad e=1+k(m-1),\quad q=e(k+1)^2,\quad c=k/2,\\
 \beta_{ab}=c+(2a-k)\delta+i(2b-k)\gamma,\qquad
 \chi(S)=\prod_{a,b=0}^k(S-\beta_{ab})^e,
 \quad E=\mathbb C[S]/(\chi),\quad A=[S].
 \end{gathered}\tag{SCA1}
\]
Here $e$ is the original nilpotence length, not an exponential. The original
divided Taylor coordinates are retained: if $x_{ab,t}$ is the coefficient
of $(S-\beta_{ab})^t$, $0\le t<e$, then
$(Ax)_{ab,t}=\beta_{ab}x_{ab,t}+x_{ab,t-1}$, where $x_{ab,-1}=0$.
Thus the following quotient calculations retain the original spectral
nilpotent. The AMC conductor flag is used on its proved simple-packet
domain $m=1$, $k\ge9$; the source and relation calculations also apply to
the complete orders displayed in SCA1.

The arithmetic measure is exactly $d\mu_{\rm ar}=w_h^{*k}(y)dy$.
For the two original Gamma orders used in the signed budget retain
\[
 dm_s(y)=\frac{(2\pi)^{s/2}2^{s/2}}{4\pi\Gamma(s/2)}
           |\Gamma(s/4+iy/2)|^2dy,\quad s\in\{1,k\},\qquad
 \omega_n^{(s)}=(2\pi)^{s/2}n!(s/2)_n.
\]
Their masses are $(2\pi)^{s/2}$ and $dm_1=d\sigma$; this is the
original MCE2/FAT source convention, including its full normalization.
The fixed Gamma measure and the original source interpolation are
\[
 d\sigma(y)=\frac{|\Gamma(1/4+iy/2)|^2}{2\pi}\,dy,
 \qquad d\mu_x=(1-x)d\sigma+x\,d\mu_{\rm ar},\quad0\le x\le1.
 \tag{SCA2}
\]
Its mass is $(1-x)\sqrt{2\pi}+x(\int w_h)^k$.
The arithmetic and Gamma densities are even and positive almost everywhere;
the full original moment estimates make all finite polynomial Grams positive.
No source is replaced by a probability measure. Unless a source superscript
or parameter is written, the following identities hold for any one of
these original measures, at its fixed parameter.

Let $\mathcal P_d=\mathbb C[S]_{\le d}$, with ordinary coefficient Gram
\[
 (M_d)_{ab}=\int\overline{(c+iy)^a}(c+iy)^b\,d\mu(y),
 \quad0\le a,b\le d.
\]
Thus $M_d$ has $d+1$ rows. Write $J_d:\mathcal P_d\to E$ for the original
remainder map in the stated primary coordinates. For $N\ge q-1$ put
\[
 G_N=(J_NM_N^{-1}J_N^*)^{-1},\qquad
 \mathfrak T_N=M_N^{-1}J_N^*G_N:E\longrightarrow\mathcal P_N.
 \tag{SCA3}
\]
The confluent evaluation of $1,S,\ldots,S^{q-1}$ is invertible, so these
inverses exist. Multiplication gives $J_N\mathfrak T_N=I$ and
$\mathfrak T_N^*M_N\mathfrak T_N=G_N$. For
$a\in\ker J_N$, one has $a^*M_N\mathfrak T_N=a^*J_N^*G_N=0$.
Consequently $\mathfrak T_N$ is the unique minimum lift and
\[
 \mathcal P_N=\chi\mathcal P_{N-q}\mathbin{\perp_{M_N}}
                      \mathfrak T_NE.
 \tag{SCA4}
\]
A negative degree means the zero space. All maps into a larger polynomial
space below use literal zero extension of their ordinary coefficient columns.

Let $p_n$ be the monic source orthogonal polynomial, $\omega_n=\|p_n\|^2$,
and $b_n=[p_n]\in E$. Their original phase is
$p_n(S)=i^nr_n((S-c)/i)$, where $r_n$ is the real monic polynomial.
In particular $\omega_0=\int d\mu$ is retained. Expansion in this complete
orthogonal frame gives
\[
 R_N=G_N^{-1}=\sum_{n=0}^N\frac{b_nb_n^*}{\omega_n},\qquad
 \mathfrak T_Nx=\sum_{n=0}^N
               \frac{p_n b_n^*G_Nx}{\omega_n}.
 \tag{SCA5}
\]
Indeed the quotient map in the orthonormal polynomial frame has columns
$b_n/\sqrt{\omega_n}$; inserting this frame into SCA3 proves both equalities.
Let $u_t$ be monic and orthogonal for the literal relation measure
$|\chi(c+iy)|^2d\mu$, and set
\[
 a_t=\chi u_t,\quad\deg a_t=q+t,\quad
 \nu_t=\|a_t\|^2>0,\qquad T_t=\nu_t/\omega_{q+t}.
 \tag{SCA6}
\]
These are the original SXR/TCN polynomials and full norms. Since $a_t$ has
the nonreal zeros of $\chi(c+iy)$, it cannot be the whole-source monic
minimizer, whose real-coordinate zeros are simple and real. Thus $T_t>1$.

\subsection{Exact adjacent section and relation formulas}

Put $\ell=N-q$ and $b_+=b_{N+1}$; when $N=q-1$, $\ell=-1$.
The exact newest monic relation and the new minimum section are
\[
 \begin{gathered}
 a_{\ell+1}=p_{N+1}-\mathfrak T_Nb_+,\qquad
 \nu_{\ell+1}=\omega_{N+1}+b_+^*G_Nb_+,\\
 \mathfrak T_{N+1}
   =\mathfrak T_N+a_{\ell+1}\frac{b_+^*G_N}{\nu_{\ell+1}},\qquad
 \mathfrak T_N^*M_{N+1}a_{\ell+1}=-G_Nb_+.
 \end{gathered}\tag{SCA7}
\]
Here and below a polynomial times a row denotes the corresponding linear
map. For proof the first difference is monic of degree $N+1$, has remainder
zero and is orthogonal to $\chi\mathcal P_{\ell}$: its first term is
orthogonal to all of $\mathcal P_N$, and the second is orthogonal to that
relation space by SCA4. These three properties identify $a_{\ell+1}$.
The two terms are orthogonal, which proves its full norm. Pairing the first
identity with $\mathfrak T_N$ proves the last identity.
The middle formula has the correct remainder and is orthogonal to every old
relation and to $a_{\ell+1}$, since the last identity gives exactly the
negative of the displayed correction. Hence it is the new minimum section.
This proof works at $N=q-1$ without any inverse at $q-2$.

Define the literal, same-cutoff relation defect
\[
 \mathfrak d_N=S\mathfrak T_N-\mathfrak T_NA:
        E\longrightarrow\chi\mathcal P_{N+1-q}\subset\mathcal P_{N+1}.
 \tag{SCA8}
\]
It is the term $\chi\mathcal Q_{\Xi,N}$ of AMC25 before composition with
the residual section; it is not the Hermitian matrix called $D_N$ in OMR.
For every $N\ge q$ its complete value is
\[
 \boxed{\mathfrak d_N
   =a_{\ell+1}\frac{b_N^*G_N}{\omega_N}
      +a_\ell\frac{b_{N+1}^*G_N}{\nu_\ell}.}
 \tag{SCA9}
\]
At the first cutoff it is exactly
\[
 \mathfrak d_{q-1}=\chi\frac{b_{q-1}^*G_{q-1}}{\omega_{q-1}}.
 \tag{SCA10}
\]
To prove SCA9, the remainder of SCA8 is zero. On the norm line
$\overline S=k-S$, so multiplication by $S$ has formal polynomial adjoint
$k-S$ whenever the two displayed polynomial degrees fit the pairing.
For $a\in\chi\mathcal P_{\ell-1}$ this gives
\[
 \langle a,\mathfrak d_Nx\rangle
 =\langle(k-S)a,\mathfrak T_Nx\rangle=0,
\]
because $(k-S)a\in\chi\mathcal P_\ell$. Thus the image is in the
orthogonal span of $a_\ell,a_{\ell+1}$. The leading coefficient of
$\mathfrak T_Nx$ is $b_N^*G_Nx/\omega_N$ by SCA5, giving the first row
in SCA9. The polynomial $Sa_\ell-a_{\ell+1}$ is an old relation of
degree at most $N$.
Therefore
\[
 \langle a_\ell,\mathfrak d_Nx\rangle
 =-\langle Sa_\ell,\mathfrak T_Nx\rangle
 =-\langle a_{\ell+1},\mathfrak T_Nx\rangle
 =b_{N+1}^*G_Nx.
\]
Division by the full positive norm $\nu_\ell$ proves the second row and
its positive sign. At $N=q-1$ the relation space in degree $q$ is spanned
by $\chi$ alone, so the leading-coefficient argument proves SCA10.

In particular AMC25's previously implicit ordinary quotient is now
\[
 \mathcal Q_{\Xi,N}
  =\left(u_{\ell+1}\frac{b_N^*G_N}{\omega_N}
        +u_\ell\frac{b_{N+1}^*G_N}{\nu_\ell}\right)S_\Xi
 \quad(N\ge q),
 \tag{SCA11}
\]
and is $b_{q-1}^*G_{q-1}S_\Xi/\omega_{q-1}$ at the first cutoff.
Its degree is at most $N+1-q$, and its image has rank at most two
(at most one at the first cutoff). No chosen relation coefficient is omitted.

\subsection{The full norm and the formerly implicit cross-pairing}

Orthogonality of the two monic relation polynomials in SCA9 and the last
identity of SCA7 prove
\[
 \begin{gathered}
 \mathfrak d_N^*M_{N+1}\mathfrak d_N
 =\frac{\nu_{\ell+1}}{\omega_N^2}
      G_Nb_Nb_N^*G_N
  +\frac1{\nu_\ell}G_Nb_{N+1}b_{N+1}^*G_N,\\
 \boxed{\mathfrak T_N^*M_{N+1}\mathfrak d_N
       =-\frac1{\omega_N}G_Nb_{N+1}b_N^*G_N.}
 \end{gathered}\tag{SCA12}
\]
The second term in the first line is absent at $N=q-1$; the second line
holds there unchanged. The old relation $a_\ell$ is orthogonal to the
minimum section, but the newest relation $a_{\ell+1}$ generally is not,
since that section uses the smaller cutoff $N$. This is why the displayed
cross-pairing must be calculated rather than set to zero.

Applying $2\operatorname{Re}\langle\mathfrak T_Nx,\cdot\rangle$ to
SCA8, and using $2\operatorname{Re}S=k$ on the actual source, gives
\[
 \begin{gathered}
 A^*G_N+G_NA-kG_N
 =-\bigl(\mathfrak T_N^*M_{N+1}\mathfrak d_N
       +\mathfrak d_N^*M_{N+1}\mathfrak T_N\bigr)\\
 =\frac1{\omega_N}
    G_N(b_{N+1}b_N^*+b_Nb_{N+1}^*)G_N.
 \end{gathered}\tag{SCA13}
\]
Thus the cross-pairing is exactly the source origin of the same OMR/AMC
Hermitian action defect. It is not an unrelated positive form.

On the simple-packet AMC flag, retain all original $C,D,S_\Xi,S_C,B_\Xi$
and put $Y=\kappa+\mathcal E_C=(I-P_\Xi)AS_\Xi$ as in AMC10--12.
SCA9 and AMC26 give the complete source map
\[
 S\mathfrak T_NS_\Xi-\mathfrak T_NS_\Xi B_\Xi
             =\mathfrak T_NY+\mathfrak d_NS_\Xi.
 \tag{SCA14}
\]
If $\mathcal D_N=\mathfrak d_N^*M_{N+1}\mathfrak d_N$, its exact Gram is
\[
 \begin{aligned}
 &\bigl(S\mathfrak T_NS_\Xi-\mathfrak T_NS_\Xi B_\Xi\bigr)^*
 M_{N+1}\bigl(S\mathfrak T_NS_\Xi-\mathfrak T_NS_\Xi B_\Xi\bigr)\\
 &=\kappa^*G_N\kappa
 +(C_bS_\Xi)^*(CR_NC^*)^{-1}(C_bS_\Xi)
 +S_\Xi^*\mathcal D_NS_\Xi\\
 &\quad-\frac1{\omega_N}
 \left(Y^*G_Nb_{N+1}b_N^*G_NS_\Xi
      +S_\Xi^*G_Nb_Nb_{N+1}^*G_NY\right).
 \end{aligned}\tag{SCA15}
\]
The first two terms follow from the actual $G_N$-orthogonal kernel/conductor
decomposition in AMC11. SCA12 evaluates both cross terms, with their signs.
For any individual grade $a\in\{K,\Xi,C\}$ the same calculation gives
\[
 (\mathfrak T_NP_a)^*M_{N+1}\mathfrak d_N
   =-\frac1{\omega_N}P_a^*G_Nb_{N+1}b_N^*G_N.
 \tag{SCA16}
\]
Hence none of the mixed pairings has been removed by working in a
minimum source frame. For the fixed original section $L$ substitute
$L=S_\Xi+(L-S_\Xi)$ into SCA8--16; the map $A(L-S_\Xi)$ remains exactly
the three-part correction of AMC13. This is a linear substitution in the
proved formulas, with no invariance assertion about its kernel.

\subsection{Exact reflection and both singular values of the relation defect}

Let $\mathcal R P(S)=P(k-S)$ on polynomials and on $E$. Since $q$ is even
and the original root multiset is invariant under $\beta\mapsto k-\beta$,
$\chi(k-S)=\chi(S)$. In divided Taylor coordinates the exact formula is
\[
 (\mathcal R x)_{ab,t}=(-1)^t x_{k-a,k-b,t}.
 \tag{SCA17}
\]
Thus the original nilpotent coordinates are also transported. Evenness of
the complete measure gives a source isometry. It carries each quotient
fibre to the reflected fibre, so uniqueness of its minimum proves
$\mathcal R^*G_N\mathcal R=G_N$ and
$\mathfrak T_N\mathcal R=\mathcal R\mathfrak T_N$.
Uniqueness of monic orthogonal polynomials gives
$\mathcal Rp_n=(-1)^np_n$. Consequently
\[
 b_N^*G_Nb_{N+1}=0.
 \tag{SCA18}
\]
Both endpoint vectors are nonzero: a whole-source orthogonal polynomial
has only real zeros in $y$, whereas reduction to zero would require its
nonreal zeros to include those of $\chi$. Their $G_N$-orthogonality is
therefore an actual orthogonality of two nonzero input rows.

Set $V_N=\det G_N$ and, for $j\ge q$,
$\delta_j^{\rm ar}=V_j/V_{j-1}$ when the source is arithmetic. Write
$\delta_j$ for the identical source-defined ratio in the intermediate
calculation. The determinant lemma and SCA7 give
\[
 \delta_j=\frac{\omega_j}{\nu_{j-q}}=T_{j-q}^{-1}\in(0,1),\quad
 \frac{b_N^*G_Nb_N}{\omega_N}=1-\delta_N\quad(N\ge q),\quad
 \frac{b_{q-1}^*G_{q-1}b_{q-1}}{\omega_{q-1}}=1.
 \tag{SCA19}
\]
For the last identity the first $q$ orthogonal columns form an invertible
matrix, whose Gram inverse makes these normalized columns orthonormal.
For the middle identity use the rank-one update from $G_{N-1}$ to $G_N$;
the new normalized column's quadratic norm is its previous norm divided
by one plus that norm. This equals $1-\delta_N$.

Let $\epsilon_N$ be the original positive eigenvalue of
$G_N^{-1/2}(A^*G_N+G_NA-kG_N)G_N^{-1/2}$.
SCA13 and SCA18 prove exactly
\[
 \epsilon_N^2
   =\frac{(b_N^*G_Nb_N)(b_{N+1}^*G_Nb_{N+1})}{\omega_N^2}.
 \tag{SCA20}
\]
This is the OMR untilted quantity, with its phase term zero by the original
reflection, not a redefinition using the residual compression.
SCA9 has orthogonal output columns $a_\ell,a_{\ell+1}$ and orthogonal
input rows after conjugating the domain by $G_N^{-1/2}$. Its two squared
singular values, for the original domain metric and full source norm, are
therefore
\[
 \boxed{\frac{\epsilon_N^2}{1-\delta_{N+1}},\qquad
        \frac{\epsilon_N^2\delta_N}{1-\delta_N}}
 \quad(N\ge q).
 \tag{SCA21}
\]
At $N=q-1$ only the first value occurs. Indeed the first squared value in
SCA12 is $\nu_{\ell+1}(b_N^*G_Nb_N)/\omega_N^2$; divide SCA20 by it
and use $\nu_{\ell+1}-\omega_{N+1}=b_{N+1}^*G_Nb_{N+1}$.
For the second value use
$\nu_\ell=\omega_N/\delta_N$ and the middle identity of SCA19.
This proves every factor in SCA21. In particular the exact operator norm
at every finite cutoff is the square root of the larger displayed value;
no eventual estimate is needed to define or compute it.

\subsection{The original budget and the full observation determinant}

Here use the simple-packet AMC/OMR common domain with the actual original
period, coefficient order and certificate. Write
\[
 \mathscr S(f)=f_{q-1}+f_q-f_{2q-1}-f_{2q},\qquad
 F_E^{\mu}=\mathscr S(\log\det G_N^{\mu})=\mathcal B_k^{\mu}.
 \tag{SCA22}
\]
Let $I_K$ be the unchanged frame of $K=\ker\Lambda$ and set
$H_N=I_K^*G_NI_K$, the original \emph{kernel Gram}. It is not the signed
Gamma transition $\mathcal H_k$. On a fixed basis of the complete original
observation image $B=\Lambda(E)$ define
$Q_{B,N}=(\Lambda R_N\Lambda^*)^{-1}$.
The exact splitting determinant, in any fixed lift of that basis, gives
$F_E^{\mu}=F_K^{\mu}+F_B^{\mu}$, since its constant change-of-basis
determinant cancels under $\mathscr S$.

To identify this full $B$ with the AMC flag, put
$\mathscr O=(C,D)^{\mathsf T}:E\to\mathbb C^{q'+j}$.
It is onto: lift any $C$ value, then correct its $D$ value by a vector
in $\ker C$, on which $D$ is onto. Its kernel is the actual $K$.
Thus the exact map
\[
 \mathcal U:B\longrightarrow\mathbb C^{q'+j},\qquad
 \mathcal U(\Lambda x)=\mathscr Ox
 \tag{SCA23}
\]
is well defined, injective and onto. It retains the full observation,
including its period frame; it is not the $j$-dimensional residual alone.
Its covariance is
\[
 \mathscr OR_N\mathscr O^*
 =\begin{pmatrix}CR_NC^*&CR_ND^*\\DR_NC^*&DR_ND^*\end{pmatrix},
 \qquad
 \det(\mathscr OR_N\mathscr O^*)
       =\det(CR_NC^*)\det\Omega_N.
 \tag{SCA24}
\]
The second equality follows by triangular block elimination; its Schur
complement is exactly $D(R_N-R_NC^*(CR_NC^*)^{-1}CR_N)D^*=\Omega_N$.
Since $\mathscr O=\mathcal U\Lambda$ is fixed across the four cutoffs,
the factor $|\det\mathcal U|^2$ cancels. Therefore, exactly,
\[
 F_B^{\mu}=F_C^{\mu}+\Phi_{k,\mu},\qquad
 F_C^{\mu}=\mathscr S(-\log\det(CR_N^{\mu}C^*)),\qquad
 \Phi_{k,\mu}=\mathscr S(\log\det Q_{\Xi,N}^{\mu}).
 \tag{SCA25}
\]
Here $Q_{\Xi,N}=\Omega_N^{-1}$, with precisely its previous source metric.
The formula proves the actual coefficient/observation morphism and every
determinant factor before any cancellation is made.

Retain the original OMR13,19 definitions in full:
\[
 \begin{gathered}
 \mathcal H_k=\mathcal B_k^\sigma-\mathcal B_k^{(k)}=-\mathfrak T_k,
 \quad\delta_k^\sigma=\mathcal B_k^{\rm ar}-\mathcal B_k^\sigma,\\
 \Psi_k=F_E^{\rm ar}-F_B^{(k)}
       =F_K^{(k)}+\mathcal H_k+\delta_k^\sigma,\\
 \mathcal W_k^{\rm ar}
 =\log\frac{\omega_{2q-1}^{\rm ar}\omega_{2q}^{\rm ar}}
                 {\omega_{q-1}^{\rm ar}\omega_q^{\rm ar}},\\
 \mathcal P_-=-\log(1-\delta_{2q-1}^{\rm ar})
      -2\sum_{j=q}^{2q-2}\log(1-\delta_j^{\rm ar}),\\
 \mathcal P_+=-\log(1-\delta_q^{\rm ar})-\log(1-\delta_{2q}^{\rm ar})
      -2\sum_{j=q+1}^{2q-1}\log(1-\delta_j^{\rm ar}),\\
 C_k^{\rm budget}=F_B^{(k)}+\mathcal W_k^{\rm ar}-\mathcal P_--\mathcal P_+,
 \quad w_{q-1}=w_{2q-1}=1,
 \quad w_N=2\ (q\le N\le2q-2).
 \end{gathered}\tag{SCA26}
\]
The arithmetic phase vanishes by SCA17--18, so the original action return
and its complete cancellation are
\[
 \boxed{\begin{aligned}
 J_k^{\rm action}
 &=2\sum_{N=q-1}^{2q-1}w_N\log\epsilon_N^{\rm ar}
   =\Psi_k+C_k^{\rm budget}\\
 &=F_E^{\rm ar}+\mathcal W_k^{\rm ar}-\mathcal P_--\mathcal P_+.
 \end{aligned}}\tag{SCA27}
\]
For a direct proof, SCA19--20 give at $N=q-1$
$\epsilon_N^2=(\omega_{N+1}/\omega_N)(\delta_{N+1}^{-1}-1)$ and at
$N\ge q$
$\epsilon_N^2=(\omega_{N+1}/\omega_N)(1-\delta_N)
(\delta_{N+1}^{-1}-1)$.
Take logs with the exact weights in SCA26. The monic ratios telescope to
$\mathcal W_k^{\rm ar}$; the $-\log\delta_{N+1}$ terms give the four
cutoff determinant return; the two remaining contraction sums are precisely
$-\mathcal P_-$ and $-\mathcal P_+$. This proves SCA27 independently of
any new residual estimate. Substituting SCA25 into both occurrences of
$F_B^{(k)}$ shows its full conductor and residual pieces cancel together.
Thus a tighter kernel estimate obtained by subtracting $\Phi_{k,k}$ cannot
also hold $C_k^{\rm budget}$ fixed: its matching $+\Phi_{k,k}$ is present
in that exact budget. This is a consequence of the proved observation map,
not an assertion that the observed mixed components are irrelevant.

\subsection{Finite relation-norm return and its calculated leading coefficient}

Retain SXR14--17's exact original comparison constants
$L_{k,d}=\log(A_kD_d/a_k)$ from TCN10--11, and the explicitly computable
Gamma relation determinant
\[
 \Theta_k=\exp(-L_{k,2q})\min_{0\le t\le q}
 \frac{\det\operatorname{Gram}_\sigma(\chi S^j)_{0\le j\le t}}
      {\det\operatorname{Gram}_\sigma(\chi S^j)_{0\le j<t}}
 \frac{4^{q+t}}{\sqrt{2\pi}(2q+2t)!}.
 \tag{SCA28}
\]
The empty determinant is one. The complete proved source comparison gives
$T_t^{\rm ar}\ge\Theta_k$ simultaneously for $0\le t\le q$; its proved
original-window estimate is
$\log\Theta_k=2qa_1+O_h(k+\log q)$ with $a_1=\psi(1)>0$.
Consequently $\Theta_k\ge2$ for all sufficiently large admitted $k$ for
every fixed complete packet. No all-$k$ assumption has been imposed: SCA21
is the exact finite formula before that eventual domain as well.

On this proved eventual domain, $\delta_N\le1/2$ for every required
$N\ge q$. Hence the second squared singular value in SCA21 is at most
$\epsilon_N^2$, and its first is larger than $\epsilon_N^2$.
It follows throughout the original window that
\[
 \begin{gathered}
 \|\mathfrak d_N\|_{G_N\to M_{N+1}}^2
             =\frac{(\epsilon_N^{\rm ar})^2}{1-\delta_{N+1}^{\rm ar}},\\
 \mathfrak J_k^{\rm rel}:=
      2\sum_{N=q-1}^{2q-1}w_N\log
                  \|\mathfrak d_N\|_{G_N\to M_{N+1}}
       =J_k^{\rm action}+\mathcal P_+\\
       =F_E^{\rm ar}+\mathcal W_k^{\rm ar}-\mathcal P_-.
 \end{gathered}\tag{SCA29}
\]
This is an equality for the norm of the explicit relation operator SCA9,
not a newly named uncomputed source scalar. Before the stated eventual
domain its exact value is obtained by the maximum of the two SCA21
quantities at each index; the first cutoff still has only its first value.

For every computed $\Theta_k>1$ put
$\eta_k=-\log(1-\Theta_k^{-1})$. The original penalties obey
\[
 0\le\mathcal P_-\le(2q-1)\eta_k,\qquad
 0\le\mathcal P_+\le2q\eta_k.
 \tag{SCA30}
\]
This follows term by term from SCA19 and SCA28, with the exact number of
weights in SCA26. At $\Theta_k\ge2$, integration of $(1-t)^{-1}\le2$
for $0\le t\le\Theta_k^{-1}$ gives $\eta_k\le2\Theta_k^{-1}$.
Thus both complete contraction penalties are bounded above by
$\exp\{-2qa_1+O_h(k+\log q)\}$, with their endpoint multiplicities retained.

There is also an explicit finite evaluation of the monic source window.
In the original Gamma norm $\omega_n^\sigma=\sqrt{2\pi}(2n)!/4^n$,
so put
\[
 \mathcal W_k^\sigma
 =\log\frac{(4q-2)!(4q)!}{(2q-2)!(2q)!}-2q\log4.
 \tag{SCA31}
\]
The full TW comparison on $\mathcal P_{2q}$ gives
$\omega_n^{\rm ar}/\omega_n^\sigma\in[a_k/D_{2q},A_k]$ for all
four indices. Taking the two positive and two negative logarithms proves
\[
 |\mathcal W_k^{\rm ar}-\mathcal W_k^\sigma|\le2L_{k,2q}.
 \tag{SCA32}
\]
The common source mass cancels only after these four signed factors have
been retained. In particular $\mathcal W_k^{\rm ar}=O_h(q\log q)$,
because the factorial logarithms are $O(q\log q)$ and
$L_{k,2q}=O_h(k+\log q)$.

Combining SCA27 and SCA30--32 with the original correlated
$|\delta_k^\sigma|\le2\Xi_{k,2q}$ proves the finite bound
\[
 \left|J_k^{\rm action}-\mathcal B_k^\sigma
                         -\mathcal W_k^\sigma\right|
 \le2\Xi_{k,2q}+2L_{k,2q}+(4q-1)\eta_k
 \quad(\Theta_k>1).
 \tag{SCA33}
\]
The constant $\Xi_{k,2q}$ is the original single correlated source budget
of GTM20, not the sum of independently inflated four-endpoint bounds.
For $\mathfrak J_k^{\rm rel}$ on $\Theta_k\ge2$, the identical proof
replaces $4q-1$ by $2q-1$.

The complete original quadratic coefficient was already proved in GQB30
and TCN39:
\[
 \frac{\mathcal B_k^{\rm ar}}{q^2}\longrightarrow
 C_B=4\int_0^1\psi(t)\,dt
       =9-8\log2+\mathcal F(2,\pi)>0.
 \tag{SCA34}
\]
For clarity its actual endpoint function is specified without a new
parameter convention. With elliptic modulus $\kappa_t\in(0,1)$, write
$\mathfrak a(t)=\sqrt{1-\kappa_t^2}$ for TCN26's original endpoint ratio,
distinguishing it from the relation polynomial $a_t$ of SCA6. Thus
$t+1=E(\kappa_t)/(\mathfrak a(t)K(\kappa_t))$, and
\[
 \psi(t)=\log(1+\mathfrak a(t))+\frac t2\log(1-\mathfrak a(t)^2)
                         -(t+1)\log E(\kappa_t),
 \quad\psi(0)=\log(4/\pi),\quad\psi'(t)<0.
 \tag{SCA35}
\]
TCN26 proves $\psi(t)>0$; GQB's independent retained rational lower
certificate even gives $C_B>0.41328696166$. Its complete determinant proof
is present in the pinned original source closure. The original endpoint
argument is as follows: $F_E^{\rm ar}=\log T_0+2\sum_{t=1}^{q-1}\log T_t
+\log T_q$ by SCA19. On any fixed compact subinterval of $(0,1]$, TCN24
and its complete packet comparison give
$\log T_t=2q\psi(t/q)+O_h(k+\log q)$ uniformly. For $0\le t\le q$,
the monic trial $y^t$ and the full moment bound give
$0<\log T_t\le2(q+t)\log(4/\pi)+O_h(k+\log q)$.
Hence indices $t\le\zeta q$ contribute $O(\zeta)+o(1)$ after division
by $q^2$, and the other indices give the ordinary Riemann sum
$4\int_\zeta^1\psi$. Letting $\zeta\downarrow0$ proves the first
limit in SCA34 with every endpoint retained. This is the previous
quadratic coefficient, not a different asymptotic assigned to it.

SCA27,29--34 now evaluate the actual same-class and literal relation returns:
\[
 \boxed{\frac{J_k^{\rm action}}{q^2}\longrightarrow C_B,
 \qquad\frac{\mathfrak J_k^{\rm rel}}{q^2}\longrightarrow C_B.}
 \tag{SCA36}
\]
All limits fix the complete packet and use its original admitted orders.
This follows because $\mathcal W_k^{\rm ar}/q^2\to0$ and the two proved
penalties divided by $q^2$ vanish. In the simple-packet AMC domain
$\mathcal L_{h,k}=\delta(k+1)^3/2$, and in the full-order original OMR
domain $\mathcal L_{h,k}=\delta e(k+1)^3/2$. Both have
$\log\mathcal L_{h,k}=O_h(\log k)$. Thus the exact original geometric
mean satisfies
\[
 \frac1q\log\left(
 \frac{\prod_{N=q-1}^{2q-1}(\epsilon_N^{\rm ar})^{w_N/(2q)}}
      {\mathcal L_{h,k}}\right)\longrightarrow C_B/4>0.
 \tag{SCA37}
\]
This computes the complete aggregate allowance in the direction of a
large original action norm. It does not infer the minimum at every cutoff
from that mean, nor a bound outside the original window. The separately
proved TCN/TCX minimum estimates retain their original scope.

\subsection{Original theta class, source interpolation and observation}

For the arithmetic realization retain the full $v_h^{\otimes k}$ and
HT3's map $\mathcal V_h$, so polynomial source norms are exactly the
arithmetic norms above. On each local tensor primary the sum nilpotent
has length $e=1+k(m-1)$: its $e$th power vanishes, and its $(e-1)$st
power has nonzero coefficient
$[k(m-1)]!/[(m-1)!]^k$ on the top local monomial. The possible sum
eigenvalues are precisely the unchanged $\beta_{ab}$. Therefore
$\chi(\sum_i s_i)=0$ in the original tensor quotient by the $h(s_i)$.
Apply ordered monic division to the explicit polynomial
$\mathfrak d_Nx(\sum_i s_i)$, obtaining
\[
 \mathfrak d_Nx(\textstyle\sum_i s_i)
       =\sum_{i=1}^k h(s_i)Q_{i,N,x}(s_1,\ldots,s_k),
 \quad \deg Q_{i,N,x}\le N+1-4m.
 \tag{SCA38}
\]
The remainder has degree below $4m$ in every variable and represents zero
in that exact quotient, so is the zero polynomial. The original order of
division fixes all $Q_i$ and their earlier-variable degree restrictions.
If $h(D)F_h=\Theta\phi_h$, its complete degree-$(k-1)$ primitive is
\[
 \mathfrak P_Nx=\sum_{i=1}^k(-1)^{i-1}Q_{i,N,x}(D_1,\ldots,D_k)
 (F_h^{\otimes(i-1)}\otimes\phi_h\otimes F_h^{\otimes(k-i)}),
 \qquad d\mathfrak P_Nx=\mathcal V_h\mathfrak d_Nx.
 \tag{SCA39}
\]
The second Koszul sign from the $i-1$ preceding degree-one factors makes
each differential term positive. This is a primitive in the original full
source, or in the proved proper source $W_*$ containing these coefficients;
no claim of membership in a smaller arbitrary source is made. The full
unit is still the map identifying the cyclic polynomial quotient with its
tensor jet image. Its inverse is used only on that image.
The original $h^2$ jet and marked/proper maps are the same source restrictions
and images as TCD3--5; SCA38--39 do not replace them by unrestricted jets.
For each fixed original coefficient face and nonempty leg mask, zero-slot
insertion commutes with the polynomial coefficient maps and the ordered
primitive. Supported zero retains its face; it is not the external absolute
base $\tau$. A transported slot permutation transports that order and its
Koszul signs, rather than selecting a different primitive silently.

The SXR receiving integral is also returned to these exact sections.
At its original source parameter $x$ and $d=q+t$, SCA7 gives
\[
 a_t=p_d-\mathfrak T_{d-1}b_d,\qquad
 \mathfrak T_db_d=p_d-\frac{\omega_d}{\nu_t}a_t.
 \tag{SCA40}
\]
The second equality follows by inserting the first into the section update
and using $b_d^*G_{d-1}b_d=\nu_t-\omega_d$. Thus SXR's actual normalized
relation and quotient vectors are exactly
$g_t=a_t/\sqrt{\nu_t}$ and
$w_t=\mathfrak T_db_d/\sqrt{\omega_d(1-T_t^{-1})}$.
Its original density derivative is
$\mathcal C(x)=M_{2q}(x)^{-1}\dot M_{2q}$, distinct from the conductor $C$.
Consequently SXR11--12 give the exact additional expression
\[
 \begin{gathered}
 J_k^{\rm action}=\mathcal B_k^\sigma+\mathcal Q_k+\mathcal E_k
                  +\mathcal W_k^{\rm ar}-\mathcal P_--\mathcal P_+,\\
 \mathcal Q_k=\int_0^1\sum_{t=0}^q c_t
   \left(\langle g_t,\mathcal C(x)g_t\rangle_x
              -\langle w_t,\mathcal C(x)w_t\rangle_x\right)dx,
 \quad c_0=c_q=1,\quad c_t=2\ (0<t<q),\\
 |\mathcal E_k|\le2qL_{k,2q}\min(1,\Theta_k^{-1/2}).
 \end{gathered}\tag{SCA41}
\]
Every vector and positive phase here is identified by SCA40 with the same
section and primitive already used in SCA8--16. The source derivative is
not multiplication by $S$. For the fixed original joint kernel/image split,
the pairing $-2\operatorname{Re}\langle w_{K,t},\mathcal C(x)w_{Y,t}\rangle_x$
therefore remains the exact SXR20 term; source orthogonality alone does not
delete it. Formula SCA36 computes the $q^2$ leading action allowance without
claiming that this surviving signed integral has been evaluated to order $q$.

Finally, in the AMC simple-packet flag the source class and observation are
exactly
\[
 \Lambda\kappa=0,\qquad
 \Lambda AS_\Xi=\mathfrak b B_\Xi+\Lambda S_C C_bS_\Xi,
 \qquad \mathfrak b=\Lambda L=\Lambda S_\Xi.
 \tag{SCA42}
\]
These follow from the same original kernel and fixed section, as in AMC29.
The relation term has the explicit theta primitive SCA39; the kernel and
conductor defects retain the exact Gram SCA15 and the original observation
SCA42. In particular neither an observation-kernel class nor its action is
declared a theta boundary. The calculated norm, full determinant cancellation
and positive quadratic allowance establish these specific original-map
consequences. They prove neither RH nor an actual offcritical zero.

\subsection{The actual minimum on the whole original action window}

The preceding geometric mean is not needed to infer the minimum. The
already proved uniform original-window bound in SXR17 supplies a separate
evaluation of that minimum. Retain the exact $\Theta_k$ of SCA28 and
define $z_+=\max(z,0)$ only in the following scalar bound. Then, for
every $q-1\le N\le2q-1$, the original TW recurrence comparison and
SCA19--20 give the finite inequality
\[
 (\epsilon_N^{\rm ar})^2\ge
 q(q-\tfrac12)\exp(-L_{k,2q})
       \frac{(\Theta_k-1)_+^2}{\max(1,\Theta_k)}.
 \tag{SCA43}
\]
To prove it, put $d=N+1$. TCN13 proves
$\omega_d^{\rm ar}/\omega_{d-1}^{\rm ar}
 \ge d(d-1/2)\exp(-L_{k,d})$.
The finite constant $D_d=[1+4(d+M_h)^2]^{M_h}$ is increasing, so
$L_{k,d}\le L_{k,2q}$; also $d(d-1/2)\ge q(q-1/2)$.
When $\Theta_k>1$, both $T_{d-q}$ and, if present, $T_{d-q-1}$
are at least $\Theta_k$. The exact formula in the proof of SCA27
is therefore at least the displayed factor
$(1-\Theta_k^{-1})(\Theta_k-1)$ times this recurrence lower bound.
At $N=q-1$ its first factor is absent and is exactly one, which is
larger. If $\Theta_k\le1$, the right side of SCA43 is zero, proving
the same finite inequality without adding a cutoff assumption.
On the proved eventual domain $\Theta_k\ge2$ it is strictly positive.

For a matching upper bound on the minimum use the actual endpoint
$N=2q-1$ alone. Define the full finite Gamma reference
\[
 T_q^\sigma=
 \frac{\det\operatorname{Gram}_\sigma(\chi S^j)_{0\le j\le q}}
      {\det\operatorname{Gram}_\sigma(\chi S^j)_{0\le j<q}}
 \frac{4^{2q}}{\sqrt{2\pi}(4q)!}.
\]
The exact TCN13 upper recurrence comparison is
$\omega_{2q}^{\rm ar}/\omega_{2q-1}^{\rm ar}
 \le (2q)(2q-1/2)\exp(L_{k,2q-1})$.
Taking monic minima in both sides of the source comparison gives
$T_q^{\rm ar}\le\exp(L_{k,2q})T_q^\sigma$.
Discarding only the factors $1-T_{q-1}^{-1}<1$ and
$T_q-1<T_q$ from the exact positive formula proves
\[
 \min_{q-1\le N\le2q-1}(\epsilon_N^{\rm ar})^2
 \le (\epsilon_{2q-1}^{\rm ar})^2
 \le (2q)(2q-\tfrac12)
          \exp(L_{k,2q-1}+L_{k,2q})T_q^\sigma.
 \tag{SCA44}
\]
No estimate of a mean was used in either direction.

SXR17 proves $\log\Theta_k=2qa_1+O_h(k+\log q)$, and TCN24--26
at the fixed upper ratio $t=1$ prove
$\log T_q^\sigma=2qa_1+O_h(k+\log q)$ with the same $a_1=\psi(1)>0$.
Since $L_{k,2q}=O_h(k+\log q)$ and $\Theta_k\to\infty$, taking
logarithms in the two strictly positive eventual bounds gives
\[
 \boxed{\min_{q-1\le N\le2q-1}\log\epsilon_N^{\rm ar}
                =qa_1+O_h(k+\log q).}
 \tag{SCA45}
\]
For example the logarithm of the final factor in SCA43 is exactly
$\log\Theta_k+2\log(1-\Theta_k^{-1})$, and its second term tends
to zero. Thus no unproved uniformity near the lower ratio endpoint is
introduced: its required uniform lower bound is precisely SXR17,
whose proof includes the original small-rank interval.

Finally SCA29 bounds the logarithmic difference between the literal
relation norm and $\epsilon_N^{\rm ar}$ uniformly by
$-\tfrac12\log(1-\Theta_k^{-1})$. Taking minima preserves these
two inequalities. With the unchanged original same-class discrepancy
$\mathcal L_{h,k}=\delta e(k+1)^3/2$, this proves
\[
 \boxed{\begin{gathered}
 \lim\frac1q\min_{q-1\le N\le2q-1}
                   \log\|\mathfrak d_N\|_{G_N\to M_{N+1}}=a_1,\\
 \lim\frac1q\log\left(
    \frac{\min_{q-1\le N\le2q-1}\epsilon_N^{\rm ar}}
         {\mathcal L_{h,k}}\right)=a_1>0,\\
 \lim\frac1q\log\left(
    \frac{\min_{q-1\le N\le2q-1}
                    \|\mathfrak d_N\|_{G_N\to M_{N+1}}}
         {\mathcal L_{h,k}}\right)=a_1>0.
 \end{gathered}}\tag{SCA46}
\]
The limits fix the actual complete packet and retain exactly the original
admitted $k$ and degree window. They establish that the actual minimum
allowance throughout this window exceeds the forced polynomial discrepancy
eventually. This is stronger than the geometric-mean statement and is
proved from the uniform adjacent norm theorem, not inferred from that
mean. It rules out obtaining a vanishing loss from this particular
same-class allowance and residual subtraction; it makes no assertion
about other degree regimes or other original-program comparisons.
